\chapter{The \v{C}ech Complex}
\label{ch:cech}

The \v{C}ech complex translates the question ``do local equivalence
judgments glue to a global equivalence?'' into finite-dimensional
linear algebra.  We construct it in full, mirroring the implementation
in \texttt{cohomology.py}.

\section{Cochain Groups}

\begin{definition}[\v{C}ech Cochain Groups]
\label{def:cech-cochain}
Given a covering $\covers = \{U_i\}_{i \in I}$ of the parameter space,
the \v{C}ech cochain groups of the isomorphism sheaf are:
\begin{align}
  C^0(\covers, \Iso) &\coloneqq \prod_{i \in I} \Iso(U_i) \\
  C^1(\covers, \Iso) &\coloneqq \prod_{(i,j) \in I^{(2)}} \Iso(U_i \cap U_j) \\
  C^2(\covers, \Iso) &\coloneqq \prod_{(i,j,k) \in I^{(3)}}
    \Iso(U_i \cap U_j \cap U_k)
\end{align}
where $I^{(2)} = \{(i,j) : i < j\}$ and
$I^{(3)} = \{(i,j,k) : i < j < k\}$ are the sets of ordered pairs
and triples of covering indices.
\end{definition}

A 0-cochain $\varphi \in C^0$ is a collection of local
equivalence judgments --- one per fiber.  A 1-cochain
$\psi \in C^1$ is a collection of ``transition functions'' on
pairwise overlaps.

\begin{definition}[LocalJudgment --- Algorithmic $C^0$ Entry]
\label{def:local-judgment}
Each entry of $C^0$ is computed by the function
\texttt{\_check\_fiber} and recorded as a \emph{LocalJudgment}:
\[
  \mathsf{LocalJudgment}(U_i) \coloneqq
  \bigl(\,\mathit{fiber} : I,\;
         \mathit{is\_eq} : \{0, 1, {?}\},\;
         \mathit{cex} : \PyObj \cup \{\bot\},\;
         \mathit{concrete} : \Bool
  \bigr)
\]
Here $\mathit{is\_eq} = 1$ iff the Z3 query
$\exists x \in U_i.\;\den{f}(x) \neq \den{g}(x)$ returns
\textsc{unsat}; $\mathit{is\_eq} = 0$ iff it returns \textsc{sat}
with a concrete counterexample $\mathit{cex}$; and
$\mathit{is\_eq} = {?}$ on timeout.  The flag $\mathit{concrete}$
is $\mathsf{true}$ when the counterexample uses only interpreted
functions (no uninterpreted symbols).
\end{definition}

\begin{remark}[From LocalJudgments to $C^0$]
The 0-cochain is the vector
$\varphi = (\varphi_i)_{i \in I}$ with
$\varphi_i \coloneqq \mathit{is\_eq}(U_i) \in \mathrm{GF}(2)$
(mapping $1 \mapsto 1$, $0 \mapsto 0$, ${?} \mapsto 0$).
\end{remark}

\section{Coboundary Operators}

\begin{definition}[Coboundary $\delta^0 : C^0 \to C^1$]
\label{def:delta0}
For $\varphi = (\varphi_i)_{i \in I} \in C^0$:
\[
  (\delta^0 \varphi)_{ij} \coloneqq \restr{\varphi_j}{U_{ij}} \circ
    \restr{\varphi_i}{U_{ij}}^{-1}
\]
This measures the ``twist'' between two local equivalences on their
overlap.  If both are isomorphisms (local equivalence holds on both
fibers), the transition function is the identity.
\end{definition}

\begin{definition}[Coboundary $\delta^1 : C^1 \to C^2$]
\label{def:delta1}
For $\psi = (\psi_{ij})_{(i,j) \in I^{(2)}} \in C^1$:
\[
  (\delta^1 \psi)_{ijk} \coloneqq \restr{\psi_{jk}}{U_{ijk}} \circ
    \restr{\psi_{ij}}{U_{ijk}}^{-1} \circ \restr{\psi_{ik}}{U_{ijk}}
\]
This is the \emph{cocycle condition}: the transition functions around
every triple overlap must compose to the identity.
\end{definition}

\begin{definition}[Overlap Relation --- Implementation]
\label{def:overlap-impl}
In the implementation, two fiber combos $U_i = (t_1^i, \ldots, t_n^i)$
and $U_j = (t_1^j, \ldots, t_n^j)$ are declared to \emph{overlap}
when they share at least one axis:
\[
  U_i \cap U_j \neq \emptyset
  \quad\Longleftrightarrow\quad
  \exists\, k \in \{1, \ldots, n\}.\;
    t_k^i = t_k^j
\]
This captures the geometric fact that fibers in a product space
intersect along shared coordinate hyperplanes.
\end{definition}

\begin{lemma}[$\delta^1 \circ \delta^0 = 0$]
\label{lem:d1d0}
The composition of coboundary operators is zero:
$\delta^1(\delta^0(\varphi)) = 0$ for all $\varphi \in C^0$.
\end{lemma}

\begin{proof}
Let $\varphi = (\varphi_i)$ and set $\psi_{ij} \coloneqq (\delta^0\varphi)_{ij}
= \restr{\varphi_j}{U_{ij}} \circ \restr{\varphi_i}{U_{ij}}^{-1}$.
On a triple overlap $U_{ijk}$:
\begin{align*}
  (\delta^1\psi)_{ijk}
  &= \restr{\psi_{jk}}{U_{ijk}} \circ \restr{\psi_{ij}}{U_{ijk}}^{-1}
     \circ \restr{\psi_{ik}}{U_{ijk}} \\
  &= \bigl(\varphi_k \circ \varphi_j^{-1}\bigr) \circ
     \bigl(\varphi_j \circ \varphi_i^{-1}\bigr)^{-1} \circ
     \bigl(\varphi_k \circ \varphi_i^{-1}\bigr) \\
  &= \varphi_k \circ \varphi_j^{-1} \circ \varphi_i \circ \varphi_j^{-1}
     \circ \varphi_k \circ \varphi_i^{-1}
\end{align*}
where all restrictions are to $U_{ijk}$.  Over $\mathrm{GF}(2)$ (where
every element is its own inverse and the group is abelian), this reduces
to $0 + 0 + 0 = 0$.  In general, one may verify directly that the
six terms cancel pairwise by the group structure: $\varphi_k \varphi_k^{-1} =
\mathrm{id}$ etc., giving the identity.
\end{proof}

\section{Cohomology Groups}

\begin{definition}[\v{C}ech Cohomology Groups]
\label{def:cech-cohomology}
\begin{align}
  \Hzero(\covers, \Iso) &\coloneqq \ker(\delta^0) \\
  \Hone(\covers, \Iso) &\coloneqq \ker(\delta^1) / \mathrm{im}(\delta^0)
\end{align}
\end{definition}

\begin{remark}[Interpretation of $\Hzero$ and $\Hone$]
\label{rem:h0h1-meaning}
$\Hzero$ is the group of \emph{global sections}: 0-cochains whose local
judgments already agree on all overlaps (i.e., a consistent global
equivalence witness).

$\Hone$ measures the \emph{obstruction to gluing}.  A 1-cocycle
$\psi \in \ker(\delta^1)$ satisfies the cocycle condition on every
triple overlap, but may not arise as the coboundary of a 0-cochain.
When $\Hone \neq 0$, there exist ``twisted'' transition functions
that cannot be straightened --- the local equivalences fail to extend
globally.
\end{remark}

\begin{theorem}[The Descent Theorem]
\label{thm:descent-full}
$\Sem_f \iso \Sem_g$ (global equivalence) if and only if:
\begin{enumerate}
  \item $\Hzero(\covers, \Iso) \neq 0$ (there exist local equivalences), and
  \item $\Hone(\covers, \Iso) = 0$ (local equivalences glue to a
    global one).
\end{enumerate}
\end{theorem}

\begin{proof}
$(\Rightarrow)$ If $\Sem_f \iso \Sem_g$ globally, restricting to
each fiber gives local equivalences that trivially glue (the
transition functions are all identities).  The resulting 1-cocycle
$\psi_{ij} = \mathrm{id}$ is manifestly $\delta^0$ of the constant
global section, so $[\psi] = 0 \in \Hone$.

$(\Leftarrow)$ If local equivalences exist and $\Hone = 0$, every
1-cocycle is a coboundary: there exists $\varphi \in C^0$ with
$\delta^0 \varphi = \psi$.  In particular, the transition functions
$(\delta^0 \varphi)_{ij}$ are identities (since each $\varphi_i$ is
an isomorphism).  By the gluing axiom of the sheaf $\Iso$, the
family $(\varphi_i)$ extends uniquely to a global section
$\varphi \in \Gamma(X, \Iso)$, which is the desired global
isomorphism $\Sem_f \iso \Sem_g$.
\end{proof}

\section{Computation over GF(2)}

In our setting, the coefficient sheaf takes values in $\{0, 1\}$
(isomorphism or not).  The \v{C}ech complex reduces to a chain
complex over $\mathrm{GF}(2)$ (the field with two elements):
\[
  \mathrm{GF}(2)^{|I|} \xrightarrow{\delta^0}
  \mathrm{GF}(2)^{|I^{(2)}|} \xrightarrow{\delta^1}
  \mathrm{GF}(2)^{|I^{(3)}|}
\]

\begin{proposition}[$\Hone$ via Rank--Nullity over $\mathrm{GF}(2)$]
\label{prop:h1-rank}
Let $D_0$ and $D_1$ denote the matrix representations of $\delta^0$
and $\delta^1$ over $\mathrm{GF}(2)$.  Then:
\[
  \dim_{\mathrm{GF}(2)} \Hone
  = \dim \ker(D_1) - \dim \mathrm{im}(D_0)
  = \bigl(|I^{(2)}| - \mathrm{rank}(D_1)\bigr) - \mathrm{rank}(D_0)
\]
This is a finite linear algebra computation, decidable in
$O(n^3)$ where $n = |I^{(2)}|$ via Gaussian elimination over
$\mathrm{GF}(2)$.
\end{proposition}

\begin{proof}
By the rank--nullity theorem, $\dim \ker(D_1) = |I^{(2)}| - \mathrm{rank}(D_1)$.
The first isomorphism theorem gives
$\Hone = \ker(D_1) / \mathrm{im}(D_0)$, so
$\dim \Hone = \dim \ker(D_1) - \dim \mathrm{im}(D_0) = \dim \ker(D_1) - \mathrm{rank}(D_0)$.
(Note $\mathrm{im}(D_0) \subseteq \ker(D_1)$ by Lemma~\ref{lem:d1d0}.)
\end{proof}

\begin{definition}[$D_0$ and $D_1$ Matrices (Implementation)]
\label{def:d0d1-matrices}
Given an enumeration of fibers $U_0, \ldots, U_{n-1}$ and overlaps
$o_0, \ldots, o_{m-1}$:
\begin{itemize}
  \item $(D_0)_{k,i} = 1$ iff fiber $U_i$ participates in overlap $o_k$
    (i.e., $o_k = (U_a, U_b)$ and $i \in \{a, b\}$).
  \item $(D_1)_{t,(k)} = 1$ iff overlap $o_k$ participates in triple
    $t$ (overlaps sharing a common fiber form a triple).
\end{itemize}
The function \texttt{compute\_cech\_h1} constructs these matrices
and calls \texttt{\_gf2\_rank} (Gaussian elimination) to compute
the ranks.
\end{definition}

\begin{definition}[CechResult --- Algorithmic Output]
\label{def:cech-result}
The output of the full \v{C}ech computation is a record:
\[
  \mathsf{CechResult} \coloneqq
  \bigl(\,
    h_0 : \Nat,\;\;
    h_1 : \Nat,\;\;
    n : \Nat,\;\;
    u : \Nat,\;\;
    \mathit{obs} : \mathcal{P}(I)
  \bigr)
\]
where $h_0 = |\{i : \varphi_i = 1\}|$ (number of equivalent fibers),
$h_1 = \dim \Hone$, $n$ is the total number of fibers, $u$ is the
number of timeouts, and $\mathit{obs}$ is the set of fibers carrying
concrete counterexamples.  The verdict is:
\[
  \mathsf{CechResult}.\mathit{equivalent} =
  \begin{cases}
    \mathsf{True}  & \text{if } \mathit{obs} = \emptyset,\; h_0 > 0,\;
                      h_1 = 0,\; u = 0 \\
    \mathsf{False} & \text{if } \mathit{obs} \neq \emptyset \\
    \mathsf{None}  & \text{otherwise (inconclusive)}
  \end{cases}
\]
\end{definition}


\chapter{The Mayer--Vietoris Sequence}
\label{ch:mayer-vietoris}

The Mayer--Vietoris sequence is the principal tool for splitting a
cohomological computation into smaller, independently solvable pieces.
When a covering admits a two-piece decomposition, the long exact
sequence relates the cohomology of the whole to the cohomology of the
parts and their intersection.

\section{Setup and Statement}

\begin{definition}[Two-Piece Decomposition]
\label{def:two-piece}
Let $X$ be the parameter space with covering $\covers$.  A
\emph{two-piece decomposition} is a pair of open subsets $A, B
\subseteq X$ with $X = A \cup B$.  We write $\covers_A = \{U_i
\cap A\}$, $\covers_B = \{U_i \cap B\}$ for the induced coverings
of $A$ and $B$, and $\covers_{A \cap B}$ for the covering of
$A \cap B$.
\end{definition}

\begin{theorem}[Mayer--Vietoris]
\label{thm:mv}
For $X = A \cup B$, there is a long exact sequence:
\[
  0 \to \Hzero(X) \xrightarrow{\rho} \Hzero(A) \oplus \Hzero(B)
  \xrightarrow{\alpha} \Hzero(A \cap B)
  \xrightarrow{\delta^*} \Hone(X) \xrightarrow{\sigma} \Hone(A) \oplus \Hone(B)
  \to \Hone(A \cap B) \to \ldots
\]
where $\rho$ is the restriction map, $\alpha(s_A, s_B) \coloneqq
\restr{s_A}{A \cap B} - \restr{s_B}{A \cap B}$, and $\delta^*$
is the connecting homomorphism.
\end{theorem}

\begin{proof}[Proof sketch]
The sequence arises from the short exact sequence of cochain
complexes
\[
  0 \to C^\bullet(\covers, \Iso) \xrightarrow{\rho}
  C^\bullet(\covers_A, \Iso) \oplus C^\bullet(\covers_B, \Iso)
  \xrightarrow{\alpha}
  C^\bullet(\covers_{A\cap B}, \Iso) \to 0
\]
where $\rho$ is diagonal restriction and $\alpha$ is the difference
of restrictions.  Exactness at each level follows from the sheaf
axiom (sections that agree on an overlap patch to a section of
the union).  Applying the snake lemma / zig-zag construction to
this short exact sequence of cochain complexes yields the long
exact sequence in cohomology.
\end{proof}

\section{The Connecting Homomorphism}

\begin{definition}[Connecting Homomorphism $\delta^*$]
\label{def:connecting-hom}
Given $\gamma \in \Hzero(A \cap B)$ (a local equivalence on the
overlap), the connecting homomorphism $\delta^*(\gamma) \in \Hone(X)$
is constructed as follows:
\begin{enumerate}
  \item Since $\alpha$ is surjective on cochains, lift $\gamma$ to
    $(c_A, c_B) \in C^0(\covers_A) \oplus C^0(\covers_B)$ with
    $\alpha(c_A, c_B) = \gamma$.
  \item Apply $\delta^0$ to get $(\delta^0 c_A, \delta^0 c_B) \in
    C^1(\covers_A) \oplus C^1(\covers_B)$.
  \item The element $\delta^0 c_A$ and $\delta^0 c_B$ patch to a
    well-defined 1-cocycle on $X$ (by the exactness of the cochain
    sequence).  Its class in $\Hone(X)$ is $\delta^*(\gamma)$.
\end{enumerate}
\end{definition}

\begin{corollary}[Mayer--Vietoris Vanishing Criterion]
\label{cor:mv-vanishing}
If $\Hone(A) = \Hone(B) = 0$ and the map $\alpha : \Hzero(A) \oplus
\Hzero(B) \to \Hzero(A \cap B)$ is surjective, then $\Hone(X) = 0$.

In words: if both pieces are locally equivalent and they agree on
their overlap, then global equivalence holds.
\end{corollary}

\begin{proof}
Exactness at $\Hone(X)$ in the Mayer--Vietoris sequence gives
$\ker(\sigma) = \mathrm{im}(\delta^*)$.  Surjectivity of $\alpha$
implies $\Hzero(A \cap B) / \mathrm{im}(\alpha) = 0$, so by
exactness at $\Hzero(A \cap B)$ we have $\mathrm{im}(\delta^*) = 0$.
Exactness at $\Hone(X)$ further gives $\ker(\sigma) = 0$, and since
$\sigma$ maps into $\Hone(A) \oplus \Hone(B) = 0$, we conclude
$\Hone(X) = 0$.
\end{proof}

\section{Application: Type-Fiber Splitting}

\begin{example}[Typed vs.\ Untyped Inputs]
\label{ex:mv-typed}
A common decomposition: $A = \{$inputs where both programs agree
on types$\}$ and $B = \{$inputs where type inference is ambiguous$\}$.
If equivalence holds on both pieces and they agree on the overlap
(inputs where both type inferences apply), global equivalence follows
by Corollary~\ref{cor:mv-vanishing}.
\end{example}

\begin{proposition}[Mayer--Vietoris for the Type Fibration]
\label{prop:mv-fibration}
Let $\pi : \PyObj \to \mathsf{TypeTag}$ be the type fibration with
fibers $U_\tau = \pi^{-1}(\tau)$.  Partition the type tags into two
disjoint subsets $S_A, S_B \subseteq \mathsf{TypeTag}$ with
$S_A \cup S_B = \mathsf{TypeTag}$.  Set
$A \coloneqq \bigsqcup_{\tau \in S_A} U_\tau$ and
$B \coloneqq \bigsqcup_{\tau \in S_B} U_\tau$.  Then:
\begin{enumerate}
  \item If $S_A \cap S_B = \emptyset$ (the partition is strict), then
    $A \cap B = \emptyset$ and the Mayer--Vietoris sequence degenerates:
    $\Hone(X) \cong \Hone(A) \oplus \Hone(B)$.
  \item If the partition is not strict (some fibers are shared), the
    connecting homomorphism $\delta^*$ measures the mismatch of local
    equivalences on the shared fibers.
\end{enumerate}
\end{proposition}


\chapter{Stalks and the Stalk Criterion}
\label{ch:stalks}

We now develop the \emph{stalk} of the semantic presheaf, the
infinitesimal counterpart to the fiber-wise sections of the \v{C}ech
complex.  The stalk criterion provides the bridge between
point-by-point semantic equality and sheaf-theoretic isomorphism.

\section{Stalks of the Value Presheaf}

\begin{definition}[Stalk]
\label{def:stalk}
The \emph{stalk} of $\Sem_f$ at a point $x = (v_1, \ldots, v_n)$ is
the filtered colimit:
\[
  (\Sem_f)_x \coloneqq \varinjlim_{x \in U} \Sem_f(U) = \den{f}(v_1, \ldots, v_n)
\]
--- the program's output on the specific input $x$.  The colimit is
taken over all open sets $U$ containing $x$ in the site $\site$.
\end{definition}

\begin{definition}[Germ]
\label{def:germ}
A \emph{germ} of $\Sem_f$ at $x$ is an equivalence class $[s, U]$ where
$s \in \Sem_f(U)$ and $x \in U$.  Two germs $[s, U]$ and $[t, V]$
are equivalent if there exists $W \subseteq U \cap V$ with $x \in W$
and $\restr{s}{W} = \restr{t}{W}$.
\end{definition}

\begin{remark}
For the semantic presheaf, the stalk $(\Sem_f)_x$ is literally the
value $f(x)$: a concrete Python object.  Germs capture the idea that
two programs may be ``locally the same near $x$'' even when their
global definitions differ.
\end{remark}

\section{The Stalk Criterion}

\begin{theorem}[Stalk Criterion for Presheaf Isomorphism]
\label{thm:stalk-criterion}
Let $\eta : \presheaf \to \mathcal{G}$ be a morphism of sheaves on
a topological space (or site) $X$.  Then $\eta$ is an isomorphism
if and only if the induced map on stalks
$\eta_x : \presheaf_x \to \mathcal{G}_x$
is an isomorphism for every point $x \in X$.
\end{theorem}

\begin{proof}
$(\Rightarrow)$ Immediate: an isomorphism of sheaves induces isomorphisms
on all stalks (colimits commute with isomorphisms).

$(\Leftarrow)$ We must show $\eta$ is both injective and surjective
on sections.

\emph{Injectivity.}
Let $s, t \in \presheaf(U)$ with $\eta_U(s) = \eta_U(t)$.  For
every $x \in U$, the stalk map satisfies
$\eta_x([s, U]) = \eta_x([t, U])$.  Since $\eta_x$ is injective,
$[s, U] = [t, U]$ in $\presheaf_x$ for all $x \in U$.
That is, for each $x$ there exists $W_x \subseteq U$ with $x \in W_x$
and $\restr{s}{W_x} = \restr{t}{W_x}$.  Since $\{W_x\}_{x \in U}$
covers $U$ and $\presheaf$ is a sheaf (locality axiom), $s = t$.

\emph{Surjectivity.}
Let $r \in \mathcal{G}(U)$.  For each $x \in U$, surjectivity of
$\eta_x$ yields a germ $[s_x, V_x]$ with $\eta_x([s_x, V_x]) =
[r, U]$ in $\mathcal{G}_x$.  After shrinking $V_x$, we may
assume $\eta_{V_x}(s_x) = \restr{r}{V_x}$.  On overlaps $V_x
\cap V_y$, the sections $s_x$ and $s_y$ satisfy
$\eta(\restr{s_x}{V_x \cap V_y}) = \eta(\restr{s_y}{V_x \cap V_y})
= \restr{r}{V_x \cap V_y}$.  By the injectivity already established,
$\restr{s_x}{V_x \cap V_y} = \restr{s_y}{V_x \cap V_y}$.  The
gluing axiom of $\presheaf$ produces $s \in \presheaf(U)$ with
$\restr{s}{V_x} = s_x$, and $\eta_U(s) = r$.
\end{proof}

\begin{corollary}[Stalk Criterion for Program Equivalence]
\label{cor:stalk-equiv}
Define the morphism $\eta : \Sem_f \to \Sem_g$ by
$\eta_U(s) \coloneqq s$ whenever $\restr{\den{f}}{U} = \restr{\den{g}}{U}$.
Then $f \equiv g$ if and only if $\eta_x$ is an isomorphism for
every $x$ --- i.e., $f(x) = g(x)$ for all inputs $x$.
\end{corollary}

\begin{remark}
This is trivially true: it restates the \emph{definition} of
extensional equivalence.  The stalk criterion becomes non-trivial
when combined with the \v{C}ech complex: it justifies reducing the
check from ``all points $x$'' to ``all type fibers $U_i$'', since
each fiber is an open set in the site.
\end{remark}

\section{Symbolic Stalks}

\begin{definition}[Symbolic Stalk]
\label{def:symbolic-stalk}
Let $U_\tau = \pi^{-1}(\tau)$ be a type fiber.  The \emph{symbolic
stalk} of $\Sem_f$ on $U_\tau$ is the Z3 term
\[
  (\Sem_f)_{U_\tau} \coloneqq \den{f}\bigl|_{U_\tau}
  = \den{f}(x_1, \ldots, x_n)
  \quad\text{subject to}\quad
  \bigwedge_{k=1}^{n} \mathsf{tag}(x_k) = \tau_k
\]
where $x_1, \ldots, x_n$ are Z3 symbolic variables and
$(\tau_1, \ldots, \tau_n)$ is the fiber combo.
\end{definition}

\begin{proposition}[Symbolic Stalk Criterion]
\label{prop:symbolic-stalk}
$f \equiv g$ on fiber $U_\tau$ if and only if the Z3 query
\[
  \bigwedge_{k=1}^{n} \mathsf{tag}(x_k) = \tau_k
  \;\wedge\;
  \den{f}(x_1, \ldots, x_n) \neq \den{g}(x_1, \ldots, x_n)
\]
returns \textsc{unsat}.  This is exactly the query performed by
\texttt{\_check\_fiber} (line 626 of \texttt{checker.py}):
it asserts the fiber constraints, then asks whether the two
compiled terms can differ.
\end{proposition}

\begin{theorem}[Stalk Criterion $\Longleftrightarrow$ \v{C}ech $\Hzero$]
\label{thm:stalk-h0}
The following are equivalent:
\begin{enumerate}
  \item $\eta_x$ is an isomorphism for every $x \in U_i$ (stalk
    criterion on fiber $U_i$).
  \item $\varphi_i = 1$ in the 0-cochain (local equivalence on $U_i$).
  \item The Z3 symbolic stalk query on $U_i$ returns \textsc{unsat}.
\end{enumerate}
In particular, $\Hzero(\covers, \Iso) \neq 0$ iff there exists at
least one fiber on which the stalk criterion holds.
\end{theorem}


% ══════════════════════════════════════════════════════════
\part{The Cubical--Cohomological Synthesis}
% ══════════════════════════════════════════════════════════

\chapter{Why Neither Alone Suffices}
\label{ch:neither-alone}

Cubical type theory and \v{C}ech cohomology are each powerful on their
own, but each has a fundamental limitation when applied to program
equivalence.  Their \emph{combination} overcomes both limitations,
enabling results that neither achieves alone.  We formalize these
limitations as impossibility propositions and illustrate the failure
modes with concrete examples from the benchmark suite.

\section{What Cubical Type Theory Alone Can Do}

Cubical type theory provides:
\begin{enumerate}
  \item \textbf{Path types}: a constructive notion of equality between
    programs ($\Path_{A \to B}(f, g)$).
  \item \textbf{Function extensionality}: reduce global equivalence to
    pointwise equality ($\forall x.\, f(x) = g(x)$).
  \item \textbf{Path constructors}: equational axioms generate paths
    (commutativity, fold fusion, $\beta$-reduction, etc.).
  \item \textbf{Kan composition}: chain multiple axiom applications
    into multi-step proofs.
  \item \textbf{HITs}: quotient programs by equational theories.
\end{enumerate}

\textbf{Limitation}: cubical type theory alone treats the input space
as \emph{monolithic}.  Function extensionality says ``check $f(x) =
g(x)$ for all $x$'' --- but it gives no strategy for \emph{decomposing}
the problem when the input space is large or polymorphic.  A single
Z3 query over the full input space either succeeds (lucky --- the
programs compile to the same term) or fails (the terms are too
different for Z3 to handle).

\begin{proposition}[Cubical Incompleteness for Polymorphic Inputs]
\label{prop:cubical-incomplete}
Let $f, g : \PyObj \to \PyObj$ be two programs such that:
\begin{enumerate}
  \item $f \equiv g$ (extensionally equal on all inputs), but
  \item $\den{f} \not\equiv_{\mathrm{syn}} \den{g}$ (the compiled Z3
    terms are syntactically distinct --- they use different
    uninterpreted functions).
\end{enumerate}
Then no finite chain of equational axiom applications (cubical path
constructors) reduces $\den{f}$ to $\den{g}$ in the monolithic
setting: the Z3 query $\den{f}(x) \neq \den{g}(x)$ returns
\textsc{sat} (with a spurious model assigning uninterpreted functions
arbitrarily).
\end{proposition}

In practice: cubical path construction succeeds for 6 out of 50
equivalent pairs (the ones where the compiler produces identical Z3
terms).  The other 44 fail because the Z3 terms are structurally
different.

\section{What \v{C}ech Cohomology Alone Can Do}

\v{C}ech cohomology provides:
\begin{enumerate}
  \item \textbf{Decomposition}: break the input space into type fibers
    (a covering $\{U_i\}$).
  \item \textbf{Local checking}: verify equivalence fiber by fiber
    ($\Iso(U_i)$ for each $i$).
  \item \textbf{Gluing}: the descent theorem says local equivalences
    extend globally iff $\Hone = 0$.
  \item \textbf{Obstruction detection}: if $\Hone \neq 0$, identify
    the specific fiber where programs disagree.
\end{enumerate}

\textbf{Limitation}: cohomology alone has no mechanism for proving
that two \emph{structurally different} Z3 terms are equal on a fiber.
It can \emph{decompose} the problem into smaller queries, but each
query still asks Z3 to prove term equality.  If the terms use different
uninterpreted functions, Z3 returns SAT (false counterexample) on
every fiber.

\begin{proposition}[Cohomological Incompleteness without Normalization]
\label{prop:coh-incomplete}
There exist programs $f, g$ with $f \equiv g$ such that:
\begin{enumerate}
  \item On every fiber $U_i$, the Z3 terms $\den{f}|_{U_i}$ and
    $\den{g}|_{U_i}$ are \emph{syntactically distinct} (they use
    different uninterpreted function symbols).
  \item The \v{C}ech complex with raw Z3 terms (no normalization)
    reports $\Hone \neq 0$ on every fiber --- a false positive.
\end{enumerate}
The obstruction is not geometric but \emph{algebraic}: it arises
from the uninterpreted function symbols, not from genuine behavioral
disagreement.
\end{proposition}

In practice: the cohomological decomposition correctly identifies the
relevant fibers and correctly glues local results, but the local
checks produce the same false counterexamples as the monolithic query.

\section{Example: eq02 --- Different OTerms, $\Hone = 0$}

\begin{example}[eq02: Regex Tokenizer vs.\ State Machine]
\label{ex:eq02}
Consider the benchmark pair \texttt{eq02a} / \texttt{eq02b}:
\begin{itemize}
  \item \texttt{eq02a}: tokenizes text via \texttt{re.findall},
    counts with \texttt{Counter}, sorts by frequency.
  \item \texttt{eq02b}: tokenizes via a hand-written character-level
    state machine, counts with \texttt{defaultdict}, sorts via
    \texttt{heapq.nsmallest}.
\end{itemize}

The two implementations are extensionally equal (they produce
identical sorted word-frequency lists), but they compile to
\emph{completely different} Z3 OTerms:
\begin{itemize}
  \item $\den{\mathtt{eq02a}}$ uses uninterpreted functions for
    \texttt{re.findall} and \texttt{Counter}.
  \item $\den{\mathtt{eq02b}}$ uses uninterpreted functions for
    character predicates and \texttt{heapq.nsmallest}.
\end{itemize}

\textbf{What cubical alone sees}: the Z3 query
$\den{\mathtt{eq02a}}(x) \neq \den{\mathtt{eq02b}}(x)$ returns
\textsc{sat} immediately --- Z3 assigns uninterpreted functions to
make the terms differ.  No equational axiom closes this gap.

\textbf{What cohomology alone sees}: the covering has a single
fiber $U_{\mathsf{str}}$ (both functions take string input).  The
local check on $U_{\mathsf{str}}$ still fails for the same reason
--- the OTerms are structurally incomparable.

\textbf{What the combination sees}: on fiber $U_{\mathsf{str}}$,
the cubical normalization engine applies fold--unfold axioms and
commutativity to canonicalize both OTerms into a common normal form
(sorted word-frequency list).  With normalized terms, the Z3 query
returns \textsc{unsat}, and since there is only one fiber,
$\Hone = 0$ trivially.

This pair demonstrates that \emph{normalization within fibers}
(cubical) followed by \emph{gluing} (cohomological) succeeds
where either method alone fails.
\end{example}

\section{What the Combination Achieves}

The synthesis of cubical and cohomological methods overcomes both
limitations:

\begin{enumerate}
  \item \textbf{Cubical paths on each fiber}: the path constructors
    (equational axioms) are applied \emph{within} each fiber of the
    covering, not globally.  This is more effective because:
    \begin{itemize}
      \item On a specific fiber (e.g., $\mathsf{tag}(p) = \mathsf{Int}$),
        the Z3 terms are more constrained, and axioms like commutativity
        and fold canonicalization are more likely to close the gap.
      \item Different fibers may require \emph{different} axioms.
    \end{itemize}

  \item \textbf{Cohomological gluing of cubical paths}: the local paths
    $p_i : \Path(\den{f}|_{U_i}, \den{g}|_{U_i})$ on each fiber must
    \emph{glue} to a global path.  The \v{C}ech complex ensures
    coherence: the transition functions on overlaps must be compatible.

  \item \textbf{The cubical Kan condition as a sheaf condition}: in
    cubical type theory, the Kan condition says ``every open box has a
    filler.''  In sheaf theory, the gluing axiom says ``local sections
    that agree on overlaps extend to a global section.''  These are
    \emph{the same condition} expressed in different languages.
\end{enumerate}

\begin{theorem}[The Cubical--Cohomological Synthesis]
\label{thm:synthesis}
Let $\{U_i\}$ be a covering of the input space by type fibers.
Let $p_i : \Path(\den{f}|_{U_i}, \den{g}|_{U_i})$ be cubical paths
on each fiber, verified by Z3 with equational axioms.

Then there exists a global path $p : \Path(\den{f}, \den{g})$
if and only if $\Hone(\{U_i\}, \Iso(\Sem_f, \Sem_g)) = 0$.

Moreover, the global path is constructed by the Kan filling operation
applied to the open box formed by the local paths and the transition
functions.
\end{theorem}

\begin{proof}
\emph{(Existence $\Leftarrow$).}
Each local path $p_i$ gives a section of the path sheaf
$\Path(\Sem_f, \Sem_g)$ over $U_i$.  The transition functions on
overlaps $U_{ij}$ check that $p_i|_{U_{ij}}$ and $p_j|_{U_{ij}}$
are compatible (both prove the same equality on the overlap --- this
is automatic since they both prove $\den{f} = \den{g}$, and equality
is a proposition, so all proofs are equal by truncation).

Since $\Hone = 0$, the local sections glue to a global section by
the descent theorem.  This global section is the desired path $p$.

The Kan construction: the local paths form an open box in the product
$\prod_i \interval \to \PyObj^{U_i}$.  The Kan filling produces a
function $p : \interval \to \PyObj^X$ (over the full input space)
with $p(\mathsf{i0}) = \den{f}$ and $p(\mathsf{i1}) = \den{g}$.

\emph{(Necessity $\Rightarrow$).}
If a global path $p$ exists, restricting to each fiber gives local
paths $p|_{U_i}$.  These automatically agree on overlaps (restriction
of a global section to overlaps is coherent).  The corresponding
1-cocycle is a coboundary ($\delta^0$ of the global section), so
$\Hone = 0$.
\end{proof}

\begin{proposition}[Formalization of the Gap]
\label{prop:gap}
Define the following three decision procedures:
\begin{enumerate}
  \item $\mathsf{Cubical}(f, g)$: construct $\Path(\den{f}, \den{g})$
    via equational axioms on the monolithic input space.
  \item $\mathsf{Cech}(f, g)$: compute $\Hone$ from raw (unnormalized)
    fiber judgments.
  \item $\mathsf{CCTT}(f, g)$: normalize on each fiber (cubical), then
    compute $\Hone$ (cohomological).
\end{enumerate}
Then there exist equivalent pairs $(f, g)$ such that
$\mathsf{Cubical}(f, g) = {?}$ and $\mathsf{Cech}(f, g) = {?}$
but $\mathsf{CCTT}(f, g) = \mathsf{True}$.

Conversely, for all $(f, g)$: if $\mathsf{Cubical}(f, g) =
\mathsf{True}$ or $\mathsf{Cech}(f, g) = \mathsf{True}$, then
$\mathsf{CCTT}(f, g) = \mathsf{True}$.  The combined method
\emph{strictly extends} both individual methods.
\end{proposition}


